\newpage

% =========================================================
% ВАРИАНТ 9
% =========================================================

\section*{Вариант 9}

Сторона ромба равна \(10\), острый угол равен \(30^\circ\).
Найдите радиус окружности, вписанной в ромб.

\begin{center}
\begin{tikzpicture}[
  scale=0.62,
  line cap=round,
  line join=round,
  every node/.style={font=\small}
]
  % Ромб со стороной 6 и углом 30 градусов.
  % Масштаб рисунка не связан с числовыми данными задачи.
  \coordinate (A) at (0,0);
  \coordinate (B) at (6,0);
  \coordinate (D) at (5.196,3);
  \coordinate (C) at (11.196,3);

  % Центр вписанной окружности
  \coordinate (O) at (5.598,1.5);

  \draw[line width=0.9pt]
    (A)--(B)--(C)--(D)--cycle;

  \draw[line width=0.9pt]
    (O) circle (1.5);

  % Острый угол ромба
  \pic[
    draw,
    line width=0.65pt,
    angle radius=7mm
  ] {angle=B--A--D};

  \node at (1.35,0.38) {$30^\circ$};

  % Длина стороны
  \path (A)--(B)
    node[midway,below=3mm] {$10$};

  % Подписи
  \node[below left]  at (A) {$A$};
  \node[below right] at (B) {$B$};
  \node[above right] at (C) {$C$};
  \node[above left]  at (D) {$D$};
\end{tikzpicture}
\end{center}

\[
\boxed{\text{Ответ: }2{,}5}
\]


% =========================================================
% ВАРИАНТ 10
% =========================================================

\section*{Вариант 10}

Радиус окружности, вписанной в ромб, равен \(1{,}5\).
Найдите сторону ромба, если один из его углов равен \(30^\circ\).

\begin{center}
\begin{tikzpicture}[
  scale=0.62,
  line cap=round,
  line join=round,
  every node/.style={font=\small}
]
  \coordinate (A) at (0,0);
  \coordinate (B) at (6,0);
  \coordinate (D) at (5.196,3);
  \coordinate (C) at (11.196,3);

  \coordinate (O) at (5.598,1.5);
  \coordinate (T) at (5.598,0);

  % Ромб и вписанная окружность
  \draw[line width=0.9pt]
    (A)--(B)--(C)--(D)--cycle;

  \draw[line width=0.9pt]
    (O) circle (1.5);

  % Радиус к точке касания
  \draw[line width=0.75pt] (O)--(T);

  \pic[
    draw,
    line width=0.6pt,
    angle radius=3mm
  ] {right angle=O--T--B};

  \node[right] at ($(O)!0.5!(T)$) {$1{,}5$};

  % Угол 30 градусов
  \pic[
    draw,
    line width=0.65pt,
    angle radius=7mm
  ] {angle=B--A--D};

  \node at (1.35,0.38) {$30^\circ$};

  % Подписи
  \node[below left]  at (A) {$A$};
  \node[below right] at (B) {$B$};
  \node[above right] at (C) {$C$};
  \node[above left]  at (D) {$D$};
  \node[above left]  at (O) {$O$};
  \node[below]       at (T) {$T$};
\end{tikzpicture}
\end{center}

\[
\boxed{\text{Ответ: }6}
\]


% =========================================================
% ВАРИАНТ 11
% =========================================================

\newpage

\section*{Вариант 11}

В треугольнике \(ABC\) угол \(C\) равен \(46^\circ\),
\(AD\) и \(BE\) --- биссектрисы, пересекающиеся в точке \(O\).
Найдите угол \(AOB\). Ответ дайте в градусах.

\begin{center}
\begin{tikzpicture}[
  scale=0.68,
  line cap=round,
  line join=round,
  every node/.style={font=\small}
]
  % Треугольник с углом C, равным приблизительно 46 градусам
  \coordinate (A) at (0,0);
  \coordinate (B) at (6,0);
  \coordinate (C) at (2.5,7.037);

  % Точки пересечения биссектрис со сторонами
  \coordinate (D) at (4.441,3.135);
  \coordinate (E) at (1.082,3.047);

  % Точка пересечения биссектрис — центр вписанной окружности
  \coordinate (O) at (2.804,1.980);

  % Стороны треугольника
  \draw[line width=0.9pt]
    (A)--(B)--(C)--cycle;

  % Биссектрисы
  \draw[line width=0.85pt] (A)--(D);
  \draw[line width=0.85pt] (B)--(E);

  \fill (O) circle (1.4pt);

  % Две равные части угла A
  \pic[
    draw,
    line width=0.6pt,
    angle radius=5mm
  ] {angle=B--A--D};

  \pic[
    draw,
    line width=0.6pt,
    angle radius=5mm
  ] {angle=D--A--C};

  % Две равные части угла B
  \pic[
    draw,
    line width=0.6pt,
    angle radius=5mm
  ] {angle=C--B--E};

  \pic[
    draw,
    line width=0.6pt,
    angle radius=5mm
  ] {angle=E--B--A};

  % Угол C
  \pic[
    draw,
    line width=0.65pt,
    angle radius=6mm
  ] {angle=A--C--B};

  \node at (2.5,6.05) {$46^\circ$};

  % Подписи
  \node[below left]  at (A) {$A$};
  \node[below right] at (B) {$B$};
  \node[above]       at (C) {$C$};
  \node[right]       at (D) {$D$};
  \node[left]        at (E) {$E$};
  \node[below]       at (O) {$O$};
\end{tikzpicture}
\end{center}

\[
\boxed{\text{Ответ: }113}
\]


% =========================================================
% ВАРИАНТ 12
% =========================================================

\section*{Вариант 12}

В треугольнике \(ABC\) высота \(CH\) равна \(6\),
\(AB=BC\), \(AC=8\). Найдите синус угла \(ACB\).

\begin{center}
\begin{tikzpicture}[
  scale=0.72,
  line cap=round,
  line join=round,
  every node/.style={font=\small}
]
  % Координаты выбраны в соответствии с условием:
  % AC = 8, CH = 6, AB = BC.
  \coordinate (A) at (0,0);
  \coordinate (H) at (5.292,0);
  \coordinate (B) at (6.047,0);
  \coordinate (C) at (5.292,6);

  % Равные стороны AB и BC
  \draw[line width=0.9pt,equal segment] (A)--(B);
  \draw[line width=0.9pt,equal segment] (B)--(C);

  % Сторона AC
  \draw[line width=0.9pt] (C)--(A);

  % Высота CH
  \draw[line width=0.85pt] (C)--(H);

  % Прямой угол
  \pic[
    draw,
    line width=0.65pt,
    angle radius=3.2mm
  ] {right angle=C--H--B};

  % Числовые данные
  \path (C)--(H)
    node[midway,right=2mm] {$6$};

  \path (A)--(C)
    node[midway,above left] {$8$};

  % Подписи
  \node[below left]  at (A) {$A$};
  \node[below right] at (B) {$B$};
  \node[above]       at (C) {$C$};
  \node[below]       at (H) {$H$};
\end{tikzpicture}
\end{center}

\[
\boxed{\text{Ответ: }0{,}75}
\]

\end{document}